\subsection{表示环 R(G)}

令 \(\mathrm{R}(\mathrm{G})\) 为 \(\mathrm{G}\) 的表示环（即有限维复向量空间上的连续表示）（《代数》，第 VIII 章，第 10 节，第 6 号）。若 \(\tau\) 是 \(\mathrm{G}\) 的表示，将 \([\tau]\) 记为 \(\mathrm{R}(\mathrm{G})\) 中的等价类；若 \(\tau\) 和 \(\tau'\) 是 \(\mathrm{G}\) 的两个表示，则按定义有

\[
\begin{array}{r} [ \tau ] + [ \tau^ {\prime} ] = [ \tau \oplus \tau^ {\prime} ] \\ \relax [ \tau ] [ \tau^ {\prime} ] = [ \tau \otimes \tau^ {\prime} ]. \end{array}
\]

由于 G 的每个表示都是半单的，Z-模 \(\mathrm{R}(\mathrm{G})\) 是自由的，并以 G 的不可约表示的等价类集合为基；根据定理 1，该集合可与 \(X_{++}\) 识别。映射 \(\tau \mapsto \mathrm{L}(\tau)_{(\mathbf{C})}\) 诱导出从环 \(\mathrm{R}(\mathrm{G})\) 到环 \(\mathcal{R}(\mathfrak{g}_{\mathbf{C}})\) 的同态 l，它是 \(g_{C}\) 的表示环（第 VIII 章，第 7 节，第 6 号）。

令 \(\tau : \mathbf{G} \to \mathbf{GL}(\mathbf{V})\) 为 \(\mathbf{G}\) 的表示；考虑分次 \((\mathrm{V}_{\lambda}(\mathrm{T}))_{\lambda \in \mathrm{X}(\mathrm{T})}\)，它是 \(\mathbf{C}\)-向量空间 \(\mathbf{V}\) 的分次。将 \(\mathrm{Ch}(\mathbf{V})\) 或 \(\mathrm{Ch}(\tau)\) 记为分次向量空间 \(\mathbf{V}\) 的特征标（第 VIII 章，第 7 节，第 7 号）；若 \((e^{\lambda})_{\lambda \in \mathrm{X}(\mathrm{T})}\) 表示代数 \(\mathbf{Z}[\mathrm{X}(\mathrm{T})] = \mathbf{Z}^{(\mathrm{X}(\mathrm{T}))}\) 的典则基，则按定义有

\[
\operatorname{Ch} (\tau) = \sum_ {\lambda \in \mathrm{X} (\mathrm{T})} \left(\dim \mathrm{V} _ {\lambda} (\mathrm{T})\right) e ^ {\lambda}.
\]

以同样方式（同上引文）定义从 R(G) 到 Z[X(T)] 的环同态，仍记为 Ch。若 G 是半单的，则有交换图

\[
\begin{array}{c c c} \mathrm{R(G)} & \stackrel {{\mathrm{Ch}}} {{\longrightarrow}} & \mathbf {Z} [ \mathrm{X(T)} ] \\ \Big \downarrow_ {l} & & \Big \downarrow_ {\tilde {\delta}} \\ \mathcal {R} (\mathfrak {g} _ {\mathbf {C}}) & \stackrel {{\mathrm{ch}}} {{\longrightarrow}} & \mathbf {Z} [ \mathrm{P} ] \end{array}\tag{1}
\]

其中 P 表示 \(\mathrm{R}(\mathfrak{g}_{\mathbf{C}}, \mathfrak{t}_{\mathbf{C}})\) 的权群，\(\tilde{\delta}\) 表示由 \(\delta\) 诱导的同态。

Weyl 群 W 通过自同构作用于群 X(T)，因而作用于环 Z[X(T)]。根据第 4 节第 3 号命题 5，Ch 的像包含于子环 Z[X(T)] \(^{W}\) 中，该子环由在 W 下不变的元素组成。

命题 2. 同态 \(\mathrm{Ch}\) 诱导出从 \(\mathrm{R}(\mathrm{G})\) 到 \(\mathrm{Z}[\mathrm{X}(\mathrm{T})]^{\mathrm{W}}\) 的同构。

对于 \(\lambda \in \mathrm{X}_{++}\)，将 \([\lambda]\) 记为 \(\mathrm{R}(\mathrm{G})\) 中最高权为 \(\lambda\) 的不可约表示的等价类。由于族 \(([\lambda])_{\lambda \in \mathrm{X}_{++}}\) 是 \(\mathbf{Z}\)-模 \(\mathrm{R}(\mathrm{G})\) 的一个基，只需证明以下断言：

族 \((\mathrm{Ch}[\lambda])_{\lambda\in\mathrm{X}_{++}}\) 是 Z-模 \(\mathbf{Z}[\mathrm{X}(\mathrm{T})]^{\mathrm{W}}\) 的一个基。

对于元素 \(u = \sum_{\lambda} a_{\lambda} e^{\lambda}\)（属于 \(\mathbf{Z}[X(T)]\)），其中的项 \(a_{\lambda} e^{\lambda}\)，若 \(\lambda\) 是由 \(\mu \in X(T)\) 且 \(a_{\mu} \neq 0\) 组成的集合的最大元，则称该项为极大项。定理 1 蕴含 \(\operatorname{Ch}[\lambda]\) 有唯一的极大项，即 \(e^{\lambda}\)。因此，命题由下述引理得到。

引理 3. 对每个 \(\lambda \in \mathrm{X}_{++}\)，令 \(\mathrm{C}_{\lambda}\) 为 \(\mathbf{Z}[\mathrm{X(T)}]^{\mathrm{W}}\) 中具有唯一极大项 \(e^{\lambda}\) 的元素。则族 \((\mathrm{C}_{\lambda})_{\lambda \in \mathrm{X}_{++}}\) 是 \(\mathbf{Z}[\mathrm{X(T)}]^{\mathrm{W}}\) 的一个基。

证明与第 VI 章第 3 节第 4 号命题 3 的证明相同，只需将 A 换成 Z，将 P 换成 X(T)，并将 P ∩ \(\overline{C}\) 换成 \(X_{++}\)。

令 \(\Theta (\mathrm{G})\)（分别为 \(\Theta (\mathrm{T})\)）为取值于 \(\mathbf{C}\) 的 \(\mathrm{G}\)（分别为 T）上连续代表函数的代数，并令 \(\mathrm{Z}\Theta (\mathrm{G})\)（分别为 \(\Theta (\mathrm{T})^{\mathrm{W}}\)）为由中心函数（分别为 W-不变函数）组成的子代数。限制映射 \(\Theta (\mathrm{G})\to \Theta (\mathrm{T})\) 诱导出环同态 \(r:\mathrm{Z}\Theta (\mathrm{G})\rightarrow \Theta (\mathrm{T})^{\mathrm{W}}\)。另一方面，将表示 \(\tau\) 与其特征标对应的映射（即函数 \(g\mapsto \mathrm{Tr}\tau (g)\)）延拓为 \(\mathbf{C}\)-代数同态 \(\mathrm{Tr}:\mathbf{C}\otimes_{\mathbf{Z}}\mathrm{R}(\mathrm{G})\to \mathrm{Z}\Theta (\mathrm{G})\)；根据《谱理论》，它是同构。同样，典则嵌入 \(\mathrm{X(T)}\to \Theta (\mathrm{T})\) 诱导出 \(\mathbf{C}\)-代数同构 \(\iota :\mathbf{C}[\mathrm{X(T)}]\to \Theta (\mathrm{T})\)，进而诱导同构 \(\iota :\mathbf{C}[\mathrm{X(T)}]^{\mathrm{W}}\to \Theta (\mathrm{T})^{\mathrm{W}}\)。图

\[
\begin{array}{c c c} \mathbf {C} \otimes_ {\mathbf {Z}} \mathrm{R(G)} & \stackrel {{1 \otimes \mathrm{Ch}}} {{\longrightarrow}} & \mathbf {C} [ \mathrm{X(T)} ] ^ {\mathrm{W}} \\ \Big \downarrow_ {\mathrm{Tr}} & & \Big \downarrow_ {\iota} \\ \mathrm{Z} \Theta (\mathrm{G}) & \stackrel {{r}} {{\longrightarrow}} & \Theta (\mathrm{T}) ^ {\mathrm{W}} \end{array}\tag{2}
\]

是交换的：事实上，对于每个表示 \(\tau : \mathrm{G} \to \mathbf{GL}(\mathrm{V})\) 以及所有 \(t \in \mathrm{T}\)，

\[
\operatorname{Tr} \tau (t) = \sum_ {\lambda \in \mathrm{X} (\mathrm{T})} \left(\dim \mathrm{V} _ {\lambda} (\mathrm{T})\right) \lambda (t) = \iota (\operatorname{Ch} \tau) (t),
\]

即 \((r\circ \mathrm{Tr})[\tau ] = (\iota \circ \mathrm{Ch})[\tau ]\)

现可由命题推出以下结果。

推论. 限制映射 \(r: \mathrm{Z}\Theta(\mathrm{G}) \to \Theta(\mathrm{T})^{\mathrm{W}}\) 是双射。

